\documentclass{subfiles}
\begin{document}
    Now, we focus our attention to one of the most common problems in computer science,
        namely, the \emph{dictionary problem} which can be roughfully defined as follow:
        given a set of items \(S\), we want to define a data structure such that the following 
        three operations are performed efficiently.
        \begin{center}
            Search(S, x) \\ 
            Insert(S, x) \\ 
            Delete(S, x)
        \end{center}
    The following assumpion are considered: accessing or deleting an item \(i\) takes time \(i\),
        insertion takes time \(i + 1\) and puts the element at the end of the list.

    We begin our discussion on the problem with the following trivial remark: 
        if the elements have some kind of total order relation, then trees provide the best approach;
        if not the best possible solution is given by lists.

    Here we discuss two type of solutions, namely,
        self-adjusting lists (for the case in which no order relation exists) and self-adjusting trees
        (for the other case). Specifically, for self-adjusting trees we discuss, in order,
        red-black trees (\Cref{sec:rb-trees}) and splay trees (\Cref{sec:splay-trees}).

    \subsection{Self-adjusting lists}
    \subfile{../sub/2.1 - Self-adjusting lists}
\end{document}
